\section{FUV / ISRF: photoelectric heating and H$_2$ photodissociation}
\label{sec:fuv-isrf}

Not implemented yet -- this section records what was found while scoping
items~T5 and~T6 (Section~\ref{sec:open}): what Grackle (the frozen,
SWIFT-compatible checkout at \texttt{\textasciitilde/programs/grackle-swift/},
confirmed against its own source, not just its online docs, which are
thin on this) actually expects for these two channels
(Sections~\ref{sec:fuv-lw-grackle}--\ref{sec:fuv-two-bands}), what both
propagation designs below need in common
(Section~\ref{sec:fuv-common}), and two competing designs for turning a
per-star injected budget into a smooth ISRF: a \textbf{parabolic}
(Jacobi-relaxation) scheme (Section~\ref{sec:fuv-design-a}, 2026-07-20,
revised 2026-07-24) and a \textbf{hyperbolic} (Cattaneo/P1, upgradeable
to a cheap M1) scheme (Section~\ref{sec:fuv-design-b}, 2026-07-24,
preferred going forward because its CFL condition scales like ordinary
hydro's rather than like an explicit-diffusion $h^2$ bound). Both are
kept here, not just the latest: Design~A is simpler and may still be
the right choice if Design~B's extra machinery is not worth its cost in
practice once built, so its derivation and honest caveats are preserved
rather than deleted.

\subsection{H$_2$ photodissociation (Lyman-Werner) in Grackle}
\label{sec:fuv-lw-grackle}

Grackle has its own rate-coupling field for this, parallel to the
H/He ionization-rate fields already used for photoionization
(Section~\ref{sec:algorithm}):
\code{RT\_H2\_dissociation\_rate} (\code{grackle\_types.h:74}), threaded
through \code{solve\_chemistry.c} into the Fortran rate solver
(\code{solve\_rate\_cool\_g.F}) under the internal name
\code{kdissH2I}. Requirements, from \code{doc/source/Parameters.rst}
and \code{Integration.rst} (more explicit there than the field's own
one-line docstring):
\begin{itemize}
  \item \code{use\_radiative\_transfer=1};
  \item \code{primordial\_chemistry} (Grackle's own name for what this
    codebase calls \code{COOLING\_GRACKLE\_MODE}) set to \textbf{2 or
    3} -- i.e.\ exactly this project's modes 2/3, since a
    dissociation rate for a species (H$_2$) only means something once
    that species is tracked;
  \item units \code{1/time\_units}, the same convention as
    \code{RT\_HI\_ionization\_rate}.
\end{itemize}
Reaction \textbf{k31} in Grackle's own network is
\code{H2I + photon -> 2 HI} (\code{initialize\_rates.c:62}) -- the
Lyman-Werner dissociation reaction itself. In the solver it is applied
additively on top of the UV-background table value, exactly like
\code{kphHI} for photoionization:
\begin{quote}
\code{k31shield(i) = k31 + kdissH2I(i,j,k)} \quad
(\code{solve\_rate\_cool\_g.F:1407})
\end{quote}
\textbf{Self-shielding is already handled for us}: both
\code{H2\_self\_shielding} and \code{H2\_custom\_shielding} explicitly
attenuate \emph{both} the background UVB rate and the supplied
\code{RT\_H2\_dissociation\_rate} by the same factor -- dense,
self-shielded gas is protected from a local LW source exactly as it is
from the background, with no extra bookkeeping needed on this
codebase's side. Nothing to build in Grackle for this: the coupling
point already exists and is well-formed. The gap is entirely upstream
-- estimating the local LW flux to populate the field (see
Section~\ref{sec:fuv-band-question} below).

\subsection{Photoelectric heating in Grackle}
\label{sec:fuv-pe-grackle}

A separate, older mechanism from the RT-coupled ionization/dissociation
fields above -- not a radiative-transfer field at all, but its own
dust-physics option:

\begin{itemize}
  \item \code{photoelectric\_heating} (int, default 0): 0 = off; 1 = a
    spatially uniform rate set directly by
    \code{photoelectric\_heating\_rate} (Tasker \& Bryan 2008, no ISRF
    dependence); 2 = `Wolfire et al.\ (1995)`, Eq.~1, with a constant
    grain photoelectric efficiency $\epsilon=0.05$ (below $2\times10^4\,$K,
    0 above), scaled by the local ISRF strength; 3 = same but with
    $\epsilon$ computed directly from Wolfire et al.\ 1995 Eq.~2
    (electron-density-dependent, needs \code{primordial\_chemistry}$>0$
    to have a meaningful electron density).
  \item \code{dust\_chemistry=1} auto-enables
    \code{photoelectric\_heating=2} (if unset) \emph{and} dust
    recombination cooling (Wolfire et al.\ 1995 Eq.~9) \emph{and}
    H$_2$-formation-on-dust, all three scaled by the same ISRF value --
    i.e.\ this one flag is Grackle's intended entry point for "turn on
    dust physics with a given ISRF," not something to enable piecemeal.
  \item \textbf{The ISRF value itself can be per-particle}, not just a
    global scalar: \code{use\_isrf\_field=1} switches from the scalar
    \code{interstellar\_radiation\_field} parameter (Habing units,
    default 1.7 -- Grackle's own default happens to equal the standard
    Draine 1978 field normalization, Section~\ref{sec:fuv-band-question})
    to a per-cell/per-particle field, \code{isrf\_habing}
    (\code{grackle\_types.h:79}), same units. \textbf{This is exactly
    the plumbing item~T6 was assuming didn't exist yet} -- it does, and
    it is the more natural target for a local FUV estimate than routing
    through \code{RT\_heating\_rate} (which Section~\ref{sec:fuv-lw-grackle}'s
    ionization-rate machinery is built for, and which does not know
    about Wolfire et al.\ 1995's dust-grain physics at all).
\end{itemize}
Traced into the actual solver (\code{cool1d\_multi\_g.F:709-997}): the
per-cell ISRF value (either the scalar or, if
\code{use\_isrf\_field=1}, \code{isrf\_habing(i,j,k)}) is read into a
local \code{myisrf(i)}, and the heating rate is
\begin{equation}
  \Gamma_{\rm PE} = \gamma_{ha}\,\epsilon\,G_0\,n_H, \qquad
  \gamma_{ha} = 10^{-24}\,{\rm erg\,s^{-1}}
  \label{eq:fuv-pe-heating}
\end{equation}
(\code{rate\_functions.c}'s \code{gammah\_rate}, applied as
\code{gammaha\_eff(i) = gammaha * 0.05 * myisrf(i)} for
\code{photoelectric\_heating=2}, \code{cool1d\_multi\_g.F:972}), added
directly into \code{edot}. As with H$_2$ dissociation, nothing needs
building inside Grackle -- the per-particle \code{isrf\_habing} field
already exists; the gap is, again, estimating the local $G_0$ to fill
it.

\subsection{Concrete Grackle configuration for PE heating, dust, and a
per-particle ISRF}
\label{sec:fuv-pe-concrete}

The discussion above was in terms of what Grackle \emph{can} do;
concretely, what we would actually set (2026-07-20, not yet
implemented):

\begin{enumerate}
  \item \code{metal\_cooling=1} -- prerequisite for the next flag,
    already on for GEAR chemistry runs.
  \item \code{dust\_chemistry=1} -- Grackle's single intended entry
    point for "turn on dust physics with a given ISRF." Bundles three
    things at once, all scaled by the same ISRF value: photoelectric
    heating (auto-sets \code{photoelectric\_heating=2} if unset), dust
    recombination cooling (Wolfire et al.\ 1995 Eq.~9), and
    H$_2$-formation-on-dust (since \code{primordial\_chemistry}$>1$ in
    our modes 2/3). \textbf{Open decision, not yet made}: whether to
    take the bundle as-is, or set \code{photoelectric\_heating} and
    \code{dust\_recombination\_cooling} separately instead if
    H$_2$-formation-on-dust risks double-counting against an existing
    channel -- not checked.
  \item \code{photoelectric\_heating=2} \textbf{or} \code{=3} -- open
    decision: \code{=2} is the simpler constant-$\epsilon=0.05$
    version; \code{=3} computes $\epsilon$ from the local electron
    density (Wolfire et al.\ 1995 Eq.~2), which our
    \code{primordial\_chemistry}$>0$ tracking does provide, but
    Grackle's own docs flag that this electron density is
    "most likely underestimated at low temperatures" since it ignores
    metal/dust-grain contributions -- not yet weighed against just
    using the constant-$\epsilon$ option.
  \item \code{use\_isrf\_field=1} -- switches Grackle from the scalar
    \code{interstellar\_radiation\_field} to reading the per-particle
    \code{isrf\_habing} array (Section~\ref{sec:fuv-pe-grackle}).
  \item \textbf{Allocate and fill \code{isrf\_habing} ourselves,
    every time chemistry is solved}, one value per gas particle, in
    Habing units: take the local $G_0$-equivalent flux our own
    propagation scheme produces (physical units,
    ${\rm erg\,s^{-1}\,cm^{-2}}$, summing $L_{\rm FUV}+L_{\rm LW}$ per
    Section~\ref{sec:fuv-two-bands}) and divide by the Habing
    normalization $1.6\times10^{-3}\,{\rm erg\,s^{-1}\,cm^{-2}}$. This
    is the one number this codebase is actually responsible for
    producing; everything else in this list is Grackle configuration,
    not new code.
  \item \textbf{This is independent of H$_2$ photodissociation
    coupling} (Section~\ref{sec:fuv-lw-grackle}) -- that needs its own
    separate switches (\code{use\_radiative\_transfer=1},
    \code{primordial\_chemistry}$\in\{2,3\}$) and its own separately
    converted quantity ($L_{\rm LW}$ alone, not $L_{\rm FUV}+L_{\rm
    LW}$, into \code{RT\_H2\_dissociation\_rate}). Turning on
    \code{use\_isrf\_field} does \emph{not} also turn on H$_2$
    dissociation coupling, or vice versa -- both must be configured
    if both channels are wanted, and whether they can be enabled
    together without conflict is itself unverified
    (Section~\ref{sec:fuv-propagation}'s open-questions list).
\end{enumerate}

\subsection{Same photon band for both?}
\label{sec:fuv-band-question}

\textbf{No, not the same band.} Photoelectric heating draws on the
classical Habing band, historically $\sim6$--$13.6\,$eV
($912$--$2000$\,\AA): any FUV photon energetic enough to eject an
electron from a dust grain's surface, essentially the whole
non-ionizing far-UV continuum. Lyman-Werner (H$_2$ dissociation) is
narrower and harder: $11.2$--$13.6\,$eV ($912$--$1108$\,\AA), the
specific energy range of H$_2$'s own electronic transitions (Lyman
$B^1\Sigma_u^+$--$X^1\Sigma_g^+$ and Werner
$C^1\Pi_u$--$X^1\Sigma_g^+$ bands) -- a sub-band sitting at the hard
end of the Habing range, not the whole thing.

One escape hatch, considered but \textbf{not chosen}: assume a single
fixed ISRF spectral shape (Draine 1978 / Mathis et al.\ 1983) and
derive both band-integrated quantities from one overall FUV
normalization by a fixed ratio. Concretely, Grackle's own default
\code{interstellar\_radiation\_field=1.7} is the Habing-unit value of
exactly the standard Draine (1978) field ($\chi=1$ Draine
$\approx1.7\,G_0$ Habing) -- not a coincidence, a sign this fixed-shape
approach is that field's own default assumption -- and under that same
shape \citet{DraineBertoldi1996} give
$k_{\rm LW}\approx1.1\times10^{-10}\,\chi\,{\rm s^{-1}}$ (commonly
quoted across the PDR literature, coefficient itself not
independently re-derived from the primary text in this pass). This
would have worked, but borrows a spectral shape calibrated to the
solar neighbourhood's mixed stellar population, which does not
generally match a single star or young cluster's actual SED hardness.

\subsection{Decided (2026-07-20): two explicit sub-bands, read from
feedback tables}
\label{sec:fuv-two-bands}

Rejected the fixed-shape shortcut above in favour of tracking
\textbf{two explicit band luminosities per star}: $L_{\rm FUV}$
($6$--$11.2\,$eV) and $L_{\rm LW}$ ($11.2$--$13.6\,$eV). These are just
two more threshold-integrals of each star's own blackbody curve --
exactly the same machinery already producing $Q_H$ (threshold above
$13.6\,$eV, Section~\ref{sec:blackbody}) and $L_{\rm bol}$ (the full
integral, Section~\ref{sec:rp-momentum}), just two more threshold
pairs on the same curve. Advantages over the fixed-shape escape hatch:
\begin{itemize}
  \item Ties the FUV/LW split to each source's actual SED (which
    varies with stellar mass/temperature) instead of an imported,
    population-averaged ISM shape;
  \item For H$_2$ dissociation specifically, tracking the real
    LW-band photon flux lets the dissociation rate be built from a
    standard effective cross-section $\sigma_{\rm LW}$ (a narrow,
    well-defined atomic-physics quantity) rather than routed through
    \citeauthor{DraineBertoldi1996}'s shape-calibrated $k_{\rm
    LW}(\chi)$ fit -- one fewer inherited approximation in the chain
    (exact conversion not yet worked out -- see open items below);
  \item For photoelectric heating, sum the two bands
    ($G_0\propto L_{\rm FUV}+L_{\rm LW}$, since the Habing band is
    defined over the same $\sim6$--$13.6\,$eV range) and convert to
    Habing units via the standard normalization
    ($1\,{\rm Habing}=1.6\times10^{-3}\,{\rm erg\,s^{-1}\,cm^{-2}}$
    integrated over that range) -- direct, no shape assumption needed
    once the band integral itself is correct.
\end{itemize}
\textbf{Like $Q_H$ and $L_{\rm bol}$, both $L_{\rm FUV}$ and $L_{\rm
LW}$ will eventually be read from tables, not computed from the
current empirical piecewise power-law fits.} Item~T7
(Section~\ref{sec:open}) already plans to extend \texttt{pychem}
(\texttt{\textasciitilde/programs/pychem/}) to emit proper
$(M,Z)$-dependent stellar tables replacing today's single-metallicity
mass--luminosity/temperature fits; $L_{\rm FUV}$ and $L_{\rm LW}$ are
two more band integrals of the same underlying stellar atmosphere
model that project would need to tabulate, not a separate table-
generation effort -- widening T7's scope rather than adding a new item.

Both propagation designs below (Sections~\ref{sec:fuv-design-a}
and~\ref{sec:fuv-design-b}) consume $L_{\rm FUV}+L_{\rm LW}$ as their
source term $S$ identically; the source-term decision above is common
to both and is not repeated in either.

\subsection{Common foundations for both propagation designs}
\label{sec:fuv-common}

Material shared by Design~A (Section~\ref{sec:fuv-design-a}) and
Design~B (Section~\ref{sec:fuv-design-b}): gathered once here rather
than duplicated in both, the same principle already applied to
$\kappa_{\rm IR}$/$\kappa_{\rm FUV}$ (computed once in
Section~\ref{sec:radiation-pressure}, reused rather than re-derived).

\subsubsection{Injection alone is not a smooth ISRF}
\label{sec:fuv-propagation}

A star can inject $L_{\rm FUV}$/$L_{\rm LW}$ onto its SPH-kernel gas
neighbours using the same density-loop-style gather/scatter already
used for the HII budget (Section~\ref{sec:algorithm}) and radiation
pressure (Section~\ref{sec:radiation-pressure}). \textbf{This alone is
not a smooth ISRF}: injection only reaches particles within the
star's kernel radius $h_\star$, dropping sharply to zero outside it,
resolution-dependent (a smaller $h_\star$ at higher resolution shrinks
the field's reach, which is unphysical -- the real radiation field
does not know about the simulation's resolution), and patchy between
sources rather than the smoothly-varying background a real,
many-source, dust-attenuated ISRF produces.

\textbf{Both designs below propagate the injected budget by solving a
local, causally-bounded update instead of a full free-streaming/M1
solve at the true speed of light} (which would impose the
$\text{CFL}\sim1/c$ constraint T6's original framing was trying to
avoid), differing only in \emph{how} that bound is achieved: Design~A
(Section~\ref{sec:fuv-design-a}) solves the screened-diffusion equation
below as a quasi-steady update, relying on light-crossing time being
vastly shorter than any hydro/cooling timescale of interest; Design~B
(Section~\ref{sec:fuv-design-b}) keeps a genuinely time-dependent,
hyperbolic system but deliberately reduces its propagation speed
(a reduced-speed-of-light approximation, honestly named as one) to make
the resulting CFL condition affordable.

\subsubsection{From plain diffusion to a screened (Yukawa) solution}
\label{sec:fuv-yukawa}

Worked through in full here (2026-07-20) since it is the crux of the
propagation proposal below and easy to get wrong by analogy with
gravity/electrostatics.

\textbf{Start from ordinary diffusion.} A field $\phi$ spreading out
and fed by a source $S$ obeys
\begin{equation}
  \frac{\partial\phi}{\partial t} = D\nabla^2\phi + S .
  \label{eq:fuv-plain-diffusion}
\end{equation}
At steady state ($\partial\phi/\partial t=0$) this is just the Poisson
equation, $-D\nabla^2\phi=S$ -- the same equation gravity and
electrostatics obey. Its point-source Green's function in 3D is the
familiar
\begin{equation}
  \phi(r) \propto \frac{1}{r},
  \label{eq:fuv-poisson-green}
\end{equation}
which falls off with distance but \emph{never vanishes}: fine for
gravity (nothing removes the gravitational field), wrong for FUV
photons, which dust actually absorbs as they travel. Eq.~\ref{eq:fuv-plain-diffusion}
has no mechanism to remove $\phi$ locally, only to spread it out.

\textbf{Add a local removal (screening) term.} Physically, at every
point some of the field should just disappear -- absorbed by dust --
at a rate set by a lengthscale $\lambda$:
\begin{equation}
  -D\nabla^2\phi + \frac{\phi}{\lambda^2} = S .
  \label{eq:fuv-screened-helmholtz}
\end{equation}
This is the \textbf{screened Poisson equation}, equivalently the
\textbf{modified Helmholtz equation} -- "Helmholtz" because it has the
same form as the wave/oscillation equation $\nabla^2\phi+k^2\phi=0$,
just with the sign on the second term flipped so solutions decay
exponentially instead of oscillating. The same equation is Yukawa's
(1935) screened nuclear potential and the Debye-H\"uckel screened
potential in plasmas -- one equation, several physics contexts.

Solving Eq.~\ref{eq:fuv-screened-helmholtz} for a single point source:
away from the source ($S=0$), spherical symmetry lets the substitution
$\phi(r)=u(r)/r$ turn the 3D Laplacian into an ordinary second-order
ODE,
\begin{equation}
  u''(r) = \frac{u(r)}{D\lambda^2},
  \label{eq:fuv-radial-ode}
\end{equation}
whose physical (decaying, not growing, as $r\to\infty$) solution is
$u(r)\propto e^{-r/\ell}$ with $\ell=\sqrt{D\lambda^2}$; absorbing
constants into $\lambda$, the point-source Green's function is
\begin{equation}
  \phi(r) \;\propto\; \frac{e^{-r/\lambda}}{r},
  \label{eq:fuv-yukawa-green}
\end{equation}
the Yukawa potential: the same geometric $1/r$ spreading as
Eq.~\ref{eq:fuv-poisson-green}, now multiplied by an exponential cutoff
over the scale $\lambda$ -- exactly the shape a real photon field has
(geometric dilution \emph{and} dust extinction). Concretely,
\begin{equation}
  \lambda = \frac{1}{\kappa_{\rm FUV}(Z)\,\rho} ,
  \label{eq:fuv-lambda-def}
\end{equation}
the same "how far before dust absorbs it" quantity already entering
$\kappa_{\rm IR}\Sigma_{\rm gas}$ in the radiation-pressure term
(Eq.~\ref{eq:rp-tauIR}, Section~\ref{sec:rp-sobolev}) -- not a new
physical input, an ingredient already computed elsewhere in this
codebase, with $\kappa_{\rm FUV}=2000\,(Z/Z_\odot)\,{\rm cm^2\,g^{-1}}$
from the same \citet{Hopkins2020} fit as $\kappa_{\rm IR}$
(Section~\ref{sec:radiation-pressure}).

\subsubsection{Interface consistency: symmetrize, don't average
independently}
\label{sec:fuv-harmonic-mean}

A general problem for both designs, not just Design~A's: any local
quantity built from $\kappa_{\rm FUV}(Z_i)\rho_i$ (a decay rate, a
relaxation time, a corrected gradient -- whatever the specific design
below computes locally per particle) can differ sharply between two
neighbours at a density interface -- e.g.\ a dense, high-opacity clump
particle immediately next to a diffuse, low-opacity particle. Letting
each particle independently use its own value for a quantity that is
really a property of the \emph{link} between them (not of either
particle alone) risks an inconsistent, order-dependent exchange: it
matters whether the "resistance" of a link is dominated by the more
opaque side (physically correct -- a photon crossing the link must
still get past the denser medium) or effectively averaged away
(physically wrong -- lets a diffuse-side particle under-count how much
a neighbouring dense clump actually blocks).

\textbf{General fix, following the standard SPH treatment of
discontinuous conductivity} \citep{ClearyMonaghan1999}: combine the two
particles' local rates with a \textbf{harmonic mean}, not an
arithmetic one, the same series-resistance argument used for heat
conduction across a material interface,
\begin{equation}
  \kappa_{ij} = \frac{2\,\kappa_{{\rm FUV},i}\,\kappa_{{\rm FUV},j}}
                     {\kappa_{{\rm FUV},i}+\kappa_{{\rm FUV},j}},
  \label{eq:fuv-harmonic-mean}
\end{equation}
so the link is dominated by (but not capped exactly at) whichever side
attenuates more strongly. Design~A applies this directly to build a
single per-link relaxation rate (Section~\ref{sec:fuv-adaptive-alpha});
Design~B needs the same discipline applied to its own per-particle
correction matrices (Section~\ref{sec:fuv-b-gradients}). Purely local
quantities (e.g.\ the decay/absorption rate acting \emph{within} a
single particle's own kernel volume, not shared with a neighbour) stay
per-particle and unsymmetrized -- only link-shared quantities need
this treatment. (\citet{ClearyMonaghan1999} recalled from general SPH
heat-conduction literature -- exact volume/page not independently
re-verified this pass, unlike the primary-source checks elsewhere in
this logbook; confirm before this is cited outside an internal note.)

\subsubsection{Composability with the metagalactic/interstellar UV
background}
\label{sec:fuv-background-composability}

Neither design below produces anything but the \emph{local,
stellar-sourced} contribution to $G_0$. Neither is meant to reproduce a
metagalactic or galaxy-averaged background: the codebase already has a
separate, spatially-uniform, time-varying UV background at the Grackle
level (\code{GrackleCooling:with\_UV\_background}, present in every
current example's \code{params.yml}), and a spatially-varying
background (from the collaboration's SPHINX simulation data) is
separately planned. \textbf{Either design composes additively with
both, and duplicates neither}: the per-particle \code{isrf\_habing}
field (Section~\ref{sec:fuv-pe-grackle}) is meant to hold \emph{the
local $u_i$ alone}, converted to Habing units and \emph{added} to
whatever the background module (uniform-in-time today, SPHINX-based
map eventually) supplies at that particle's position -- not a
replacement for it. No background floor term needs adding to either
design's own update equation.

\subsection{Design A: parabolic Jacobi relaxation (2026-07-20, revised
2026-07-24)}
\label{sec:fuv-design-a}

\subsubsection{A concrete discretization: SPH smoothing with decay and
injection}
\label{sec:fuv-discretization}

The naive starting point for actually implementing
Section~\ref{sec:fuv-propagation} is SPH smoothing, the same
neighbour-loop machinery already used for metal mixing. A concrete
per-particle update rule (proposed 2026-07-20, not yet built):
\begin{equation}
  u_i^{n+1} = (1-\alpha)\,u_i^n\,e^{-\Delta t/\Delta t_{\rm decay}}
    + \alpha\sum_j \tilde w_{ij}\,u_j^n
    + \frac{u_{{\rm inject},i}}{m_i} ,
  \label{eq:fuv-discrete-update}
\end{equation}
where $u_i$ is particle $i$'s local FUV/LW field value, $\tilde w_{ij}$
a normalized SPH kernel weight, $\alpha$ a mixing fraction, and
$u_{{\rm inject},i}$ the budget gathered this step from any star whose
kernel reaches particle $i$ (Section~\ref{sec:fuv-propagation}).

\textbf{This is not an ad hoc smoothing -- it is an explicit
(forward-Euler) integrator of Eq.~\ref{eq:fuv-screened-helmholtz}.}
Expanding both exponentials to first order in $\Delta t$ and
dividing through by $\Delta t$:
\begin{equation}
  \frac{u_i^{n+1}-u_i^n}{\Delta t} \approx
    \underbrace{\frac{\alpha}{\Delta t}\Big(\sum_j\tilde w_{ij}u_j^n - u_i^n\Big)}_{\textstyle D\nabla^2 u}
    \;\underbrace{-\;\frac{u_i^n}{\Delta t_{\rm decay}}}_{\textstyle -\,u/\lambda^2}
    \;+\; \frac{u_{{\rm inject},i}}{m_i\Delta t} ,
  \label{eq:fuv-discrete-continuum-limit}
\end{equation}
which is exactly $\partial u/\partial t = D\nabla^2 u - u/\lambda^2 + S$
(Eqs.~\ref{eq:fuv-screened-helmholtz}--\ref{eq:fuv-radial-ode}), with
the correspondence
\begin{equation}
  D \sim \frac{\alpha h^2}{\Delta t}, \qquad
  \lambda^2 = D\,\Delta t_{\rm decay} .
  \label{eq:fuv-discrete-correspondence}
\end{equation}
\textbf{Consequence: $\Delta t_{\rm decay}$ is not a free tuning knob.}
Once $\lambda$ is fixed by the physical dust opacity
(Eq.~\ref{eq:fuv-lambda-def}) and $D$ by whatever $\alpha$/$h$/$\Delta t$
the smoothing pass uses, $\Delta t_{\rm decay}=\lambda^2/D$ is
determined, not chosen by feel.

\textbf{Which existing GEAR chemistry backend is the right template --
neither, exactly, but one is much closer.} SWIFT has two: plain
\code{src/chemistry/GEAR/} is a one-sided kernel-weighted
\emph{average} of the metal mass fraction (\code{chemistry\_iact.h:45-108},
\code{chemistry\_end\_density}) -- not an exchange of the metal
reservoirs themselves, and explicitly non-conservative (its own
\code{chemistry\_timestep()} returns \code{FLT\_MAX} with the comment
"no diffusion"). \code{src/chemistry/GEAR\_DIFFUSION/}
(\code{chemistry\_iact.h:169-266}, Shen et al.\ 2010) is real
metal-mass diffusion: a \emph{symmetric pairwise exchange} between
particles $i,j$ (\code{metal\_mass\_dt[i] += coef\_i*dm};
\code{metal\_mass\_dt[j] -= coef\_j*dm}), conservative to machine
precision because total metal mass must never be created or destroyed
by mixing. \textbf{Neither has anything like the $e^{-\Delta
t/\Delta t_{\rm decay}}$ term} -- both are pure redistribution.
Since our field \emph{should} lose budget locally (dust actually
absorbs FUV photons -- the sink is a feature, not a bug, unlike metal
mass), the non-conservative plain-\code{GEAR} kernel-average is the
closer starting template, not the conservative Shen et al.\ 2010
scheme built specifically to prevent exactly the kind of loss we want
here. The decay and injection terms in
Eq.~\ref{eq:fuv-discrete-update} would need to be added on top either
way -- reusing the neighbour-loop/kernel machinery, not the physics.

\textbf{Stability -- worked through in full (2026-07-20), because "do
we need a CFL condition" turned out to have a sharper answer than
"yes, a mild one."}

\textit{The decay term alone needs no CFL at all.} $u_i^n\,e^{-\Delta
t/\Delta t_{\rm decay}}$ is not a forward-Euler approximation of
$du/dt=-u/\tau$ -- it is that ODE's \emph{exact analytic solution}.
$e^{-\Delta t/\tau}\in(0,1]$ for any $\Delta t\ge0$, so this half of
Eq.~\ref{eq:fuv-discrete-update} is unconditionally stable regardless
of how large $\Delta t$ is relative to $\tau$.

\textit{The neighbour-mixing term is not solved exactly, and that is
where a real bound lives.} $\alpha\sum_j\tilde w_{ij}u_j^n$ uses each
neighbour's \emph{old} (level-$n$) value -- an explicit (Jacobi) update
of the spatial coupling, not an implicit or exact one. A von Neumann
analysis (Fourier mode $u_j^n=\hat u^n e^{ik\cdot x_j}$, kernel
response $\hat W(k)\in[-W_{\min},1]$, worst case at the
shortest-wavelength mode the particle spacing resolves) gives the
growth factor
\begin{equation}
  g(k) = (1-\alpha)e^{-\Delta t/\tau} + \alpha\,\hat W(k),
  \label{eq:fuv-growth-factor}
\end{equation}
and requiring $|g|\le1$ at the worst-case mode ($\hat W\to-W_{\min}$)
gives
\begin{equation}
  \alpha \;\lesssim\; \frac{1+e^{-\Delta t/\tau}}{W_{\min}+e^{-\Delta t/\tau}} .
  \label{eq:fuv-alpha-bound}
\end{equation}
\textbf{This bound is tightest as $\Delta t\to0$} ($\approx2/(W_{\min}+1)$)
and only loosens for larger $\Delta t$ -- the decay term actively
\emph{helps} by damping the same short-wavelength modes the mixing
term risks amplifying.

\textit{Converting to a physical decision.} Via
Eq.~\ref{eq:fuv-discrete-correspondence}'s $D\sim\alpha h^2/\Delta t$,
Eq.~\ref{eq:fuv-alpha-bound} becomes the ordinary explicit-diffusion
condition $\Delta t\lesssim h^2/D$ -- real, but nowhere near the
$\text{CFL}\sim h/c$ that ruled out M1. Concretely, using this
project's own STARBENCH numbers (Section~\ref{sec:validation}):
\code{dt\_max}$\approx10^{-5}$ in internal units
$\times\,977.8\,$Myr$\,\approx10^4\,$yr, in a $10\,$pc box with
sub-pc particle spacing, so $h/c$ is of order a few \emph{years} --
$\sim3000$--$30000\times$ smaller than the hydro step. M1 would need
thousands of sub-steps per hydro step; the diffusion-type bound, even
pushed to demand faster-than-hydro-step relaxation, needs at most a
modest number of sub-cycles (tens, not thousands) to reach it.
\textbf{This $h^2/D$ bound is nonetheless the concrete weakness Design~B
(Section~\ref{sec:fuv-design-b}) is built to avoid}: it is only mild
here because $D$ is left emergent (below) rather than tied to any
target physical rate; tying it to one reintroduces exactly this
quadratic-in-$h$ cost (Section~\ref{sec:fuv-design-a-limit}).

\textbf{Decided (2026-07-20): fix $\alpha$ once, conservatively, and
never derive it from a target relaxation length.} This is the
difference between introducing a new gas timestep limiter and not.
Deriving $\Delta t$ backwards from a desired physical $D$ (a target
screening length $\lambda$, tracked as fast as possible) risks landing
on a $\Delta t$ smaller than what gas particles already use for
hydro -- exactly the kind of new limiter (like a cooling-time
constraint) this design is trying to avoid. Instead, fix a single
conservative $\alpha$ (safe for the $\Delta t\to0$ limit of
Eq.~\ref{eq:fuv-alpha-bound}, with margin) and run the update at
whatever $\Delta t$ a gas particle \emph{already} has from its own
existing hydro/Courant timestep -- since the bound only loosens for
larger $\Delta t$, a fixed $\alpha$ chosen safe at the tightest case is
provably safe for every particle's actual, possibly much larger,
$\Delta t$. \textbf{Consequence}: no new \code{\_timestep()} criterion,
no new limiter, nothing added to the timestep hierarchy at all --
piggyback the update directly onto an already-existing gas-gas
neighbour loop (density loop, or the plain \code{GEAR/chemistry.h}
kernel-smoothing pass if active) rather than standing up a new task
subtype, since per this project's own task-registration checklist,
that is real multi-file engineering cost independent of how simple the
underlying math is.

\textbf{Trade accepted}: with $\alpha$ fixed rather than derived, the
screening length $\lambda$ (Section~\ref{sec:fuv-yukawa}) is no longer
a precisely dialled-in physical target -- it becomes an
\emph{emergent} quantity varying with local $h$ and each particle's
own $\Delta t$. Given the priority of avoiding a new gas timestep
constraint, that is the right trade: verify \emph{after the fact},
diagnostically, that the emergent relaxation time stays fast against
stellar-lifetime timescales (very likely true, since Courant-limited
hydro steps in resolved star-forming gas are already far below a Myr),
rather than enforcing an exact $\lambda$ upfront at the cost of a new
limiter. Also don't fire the update every hydro step by default -- reuse
the adaptive-rebuild-cadence machinery already built for the HII budget
(T4) so the smoothing pass runs only as often as source evolution
(stellar lifetimes, $\gtrsim$Myr) actually requires.

\subsubsection{Per-particle adaptive relaxation}
\label{sec:fuv-adaptive-alpha}

Reviewed 2026-07-24 against a design conversation (with Gemini)
re-examining whether $\alpha$ in Eq.~\ref{eq:fuv-discrete-update} should
stay a single global constant (the decision above) or vary per particle.
Conclusion: \textbf{make it per-particle} -- this is both safe and
strictly better than the fixed choice, using only the stability bound
already derived above (Eq.~\ref{eq:fuv-alpha-bound}), not a new physical
assumption.

\textbf{The bound is already local.} Eq.~\ref{eq:fuv-alpha-bound} was
derived per Fourier mode, but nothing in the derivation assumed a single
global $\Delta t$ or $\tau_{\rm decay}$ -- each particle $i$ has its own
hydro/Courant step $\Delta t_i$ and its own local decay time
$\tau_i=\lambda_i^2/D$ (set by its own local dust opacity
$\kappa_{\rm FUV}(Z_i)\rho_i$, Eq.~\ref{eq:fuv-discrete-correspondence}).
Evaluating the same bound with particle $i$'s own values,
\begin{equation}
  \alpha_i \;\lesssim\; \min\!\left(1,\;
    \frac{1+e^{-\Delta t_i/\tau_i}}{W_{\min}+e^{-\Delta t_i/\tau_i}}\right),
  \label{eq:fuv-alpha-bound-i}
\end{equation}
gives a per-particle ceiling that is provably safe for that particle's
actual $\Delta t_i$, and -- since the bound only loosens as $\Delta t_i$
grows -- is never more restrictive than a single global $\alpha$ chosen
safe for the worst (smallest $\Delta t_i/\tau_i$) particle in the whole
run. Using Eq.~\ref{eq:fuv-alpha-bound-i} in place of a fixed $\alpha$
therefore relaxes faster wherever it safely can (particles whose own
hydro step is long compared to their local decay time), with zero loss
of the stability margin the fixed choice already relied on for every
other particle. \textbf{No new timestep constraint is introduced}:
$\Delta t_i$ is read from the step the particle already has, not solved
for -- the same design principle as the original fixed-$\alpha$
decision, just applied per particle instead of globally. This
supersedes the "fix $\alpha$ once" choice above -- not because that
choice was wrong (it remains a correct conservative bound), but because
the same bound admits a strictly better per-particle evaluation once
actually written out.

\textbf{Interface treatment.} Applying the general symmetrization
principle of Section~\ref{sec:fuv-harmonic-mean} concretely here: for
the specific $i$--$j$ link, use the harmonic mean of the two
particles' local absorption rates to set the pairwise coupling
strength that enters \emph{both} particles' updates for that link,
Eq.~\ref{eq:fuv-harmonic-mean}, and symmetrize the kernel itself,
$\bar W_{ij}=\tfrac12\big(W(r_{ij},h_i)+W(r_{ij},h_j)\big)$, the same
averaging already standard in SPH force loops with differing smoothing
lengths. The purely local decay term ($\gamma_i=1/\tau_i$, dust
absorption within particle $i$'s own kernel volume) stays per-particle
and unsymmetrized -- it is not a shared link property, only the
pairwise exchange is.

\subsubsection{What Gemini's design got right, and what it over-claimed}
\label{sec:fuv-gemini-review}

A conversation exploring this same idea externally (with Gemini,
transcript reviewed 2026-07-24) reached the same practical conclusion --
per-particle adaptive relaxation is worth having -- but wrapped it in
framing that does not actually describe the scheme above, plus one
claim that is internally inconsistent. Recorded here so the corrected
framing does not get re-lost the next time this design is revisited.

\begin{itemize}
  \item \textbf{Not a reduced-speed-of-light (RSLA) scheme.} Gemini
    cast $\alpha_i$ as a per-particle $c_{{\rm eff},i}$ in a hyperbolic
    telegrapher's equation, $\frac{1}{c_{\rm eff}^2}\partial_t^2u+\ldots$.
    Eq.~\ref{eq:fuv-discrete-update} has no second time-derivative and
    no speed-of-light anywhere in it -- it is a first-order-in-time
    explicit (Jacobi) update of a parabolic screened-diffusion equation.
    What Eq.~\ref{eq:fuv-alpha-bound-i} actually is: the already-derived
    linear-stability ceiling on the mixing fraction, evaluated with each
    particle's own $\Delta t_i$ instead of the worst case in the
    simulation. Same practical behaviour (particles with locally short
    decay times relax faster), different and much simpler
    justification -- no new differential equation, no reduced-light-speed
    analogy needed \emph{for this design}. (Design~B,
    Section~\ref{sec:fuv-design-b}, deliberately \emph{does} adopt a
    genuine hyperbolic system with an honestly-named reduced propagation
    speed -- the point here is only that Design~A's own adaptive-$\alpha$
    fix does not need that machinery to work.)
  \item \textbf{The Eddington-closure/RTE derivation is not needed, and
    is internally inconsistent in the one regime it was invoked for.}
    Gemini derived $D=\tfrac13c\lambda_{\rm mfp}$ from a moment-closure
    of the radiative transfer equation, valid only when
    $\lambda_{\rm mfp}\ll h$ (optically thick -- frequent scattering
    randomizes photon direction locally). It then applied this same
    formula to the \emph{pristine-gas} limit ($\lambda_{\rm
    mfp}\to\infty$) to argue $\tau_{\rm diff}\to0$ and the scheme
    "resolves" cleanly into free streaming. That argument uses the
    optically-thick diffusion-coefficient formula precisely where it is
    not valid -- $\lambda_{\rm mfp}\gg h$ is the regime the formula
    assumes cannot happen. Design~A never needed this derivation:
    $D$ is emergent from the numerics ($D\sim\alpha h^2/\Delta t$,
    Eq.~\ref{eq:fuv-discrete-correspondence}), not from a physical mean
    free path, and the only physical input is the screening length
    $\lambda$ (Eq.~\ref{eq:fuv-lambda-def}), entering solely through the
    local decay term. Section~\ref{sec:fuv-propagation}'s original
    justification -- that the screened form is chosen because its
    Green's function \emph{shape} (Yukawa) matches real
    geometric-dilution-plus-extinction, not because direct FUV transport
    is itself a diffusion process -- was already the honest version of
    this argument and needs no revision. (Design~B below \emph{does}
    derive a genuine Eddington-closure system properly, from the actual
    moment equations, with its regime of validity stated explicitly --
    resolving this exact critique rather than repeating Gemini's
    mistake.)
  \item \textbf{The pristine-gas limit is not literally free streaming,
    and should be verified, not assumed.} At $\alpha_i\to1$ the update
    becomes a pure kernel-weighted average of neighbours each step
    (Eq.~\ref{eq:fuv-discrete-update} with no memory term) -- this can
    only move energy one kernel radius per timestep, an effective speed
    $\sim h_i/\Delta t_i$, exactly as Gemini itself computed, then
    undersold as "$\sim0.017c$, close enough to instantaneous." For a
    well-resolved star-forming region (small $h_i$, small
    Courant-limited $\Delta t_i$), this ratio is not guaranteed to stay
    fast compared to the timescales the field is meant to track (stellar
    lifetimes, cooling times) -- it is resolution- and
    timestep-dependent, and the diagnostic-verification open item below
    already covers exactly this. Treat it as something to check per
    configuration, not something the algebra guarantees.
\end{itemize}

\subsubsection{Design~A's own cost ceiling: the $h^2$ constraint}
\label{sec:fuv-design-a-limit}

Even with per-particle adaptive $\alpha_i$, Design~A's mixing term
remains a parabolic, explicit-diffusion update: its stability bound
(Eq.~\ref{eq:fuv-alpha-bound-i}), read through the correspondence
$D\sim\alpha h^2/\Delta t$, is fundamentally an
$\Delta t\lesssim h^2/D$ constraint. As implemented ($D$ left emergent,
never tied to a target physical rate), this cost is dodged only by
\emph{not} controlling the propagation speed at all -- which is also
why the emergent relaxation speed $h_i/\Delta t_i$
(Section~\ref{sec:fuv-gemini-review}, third point) is an
after-the-fact diagnostic rather than a designed, guaranteed quantity.
The moment $D$ is deliberately tied to a target propagation speed (e.g.\
to guarantee the field relaxes faster than some multiple of the sound
speed, rather than merely hoping it does), the $h^2$ scaling becomes the
concrete, quantitatively prohibitive cost Darwin flagged
(2026-07-24): at fixed target speed, halving $h$ (doubling resolution)
quarters the allowed $\Delta t$, whereas the hydro Courant condition
only halves it. \textbf{This is exactly why Design~B
(Section~\ref{sec:fuv-design-b}) exists}: a hyperbolic reformulation of
the same physics whose CFL condition scales linearly in $h$, like
ordinary hydro, letting the propagation speed be a deliberately chosen,
verified quantity rather than an emergent, hoped-for one.

\subsubsection{Open items, Design A}
\label{sec:fuv-design-a-open}

\begin{itemize}
  \item Which existing gas-gas neighbour loop to piggyback the update
    on (density loop vs.\ the plain \code{GEAR/chemistry.h}
    kernel-smoothing pass, if active in this configuration) -- not
    checked;
  \item per-particle $\alpha_i$ (Eq.~\ref{eq:fuv-alpha-bound-i}) and the
    harmonic-mean pairwise coupling (Eq.~\ref{eq:fuv-harmonic-mean}) are
    decided but not yet implemented or tested;
  \item diagnostically confirming, once built, that the emergent
    relaxation speed ($h_i/\Delta t_i$, Section~\ref{sec:fuv-gemini-review})
    stays fast against stellar-lifetime and cooling timescales, rather
    than assuming it -- not done;
  \item the propagation scheme's normalized weights $\tilde w_{ij}$ and
    how they should be built from the existing SPH kernel
    infrastructure -- not scoped.
\end{itemize}

\subsection{Design B: hyperbolic P1 relaxation, upgradeable to a cheap
M1 (2026-07-24)}
\label{sec:fuv-design-b}

Motivated directly by Section~\ref{sec:fuv-design-a-limit}'s $h^2$
ceiling: Darwin's own \texttt{swiftsim\_0} checkout already implements
a working template for exactly this kind of upgrade -- a hyperbolic
(Cattaneo/Maxwell-type) reformulation of parabolic metal diffusion,
\code{src/chemistry/GEAR\_FVPM\_DIFFUSION/hyperbolic/}, documented in
\texttt{theory/ParabolicDiffusion/GEAR\_FVPM\_Diffusion\_Theory.tex}.
That module's own timestep section states the identical trade-off in
its own words: the parabolic mode's CFL is
$\Delta t\le C_{\rm CFL}\,\Delta x^2/\|K\|_F$ (quadratic in particle size),
while its hyperbolic sibling uses an ordinary hyperbolic condition
$\Delta t = C_{\rm CFL}\,\Delta x/v_{\max}$ (linear, like hydro), with the
parabolic equation recovered exactly as the relaxation-time-to-zero
limit of the hyperbolic one. This section works out the same idea for
the FUV/ISRF field, derived from the radiative transfer equation
itself (not asserted by analogy), and identifies exactly which of its
own existing SWIFT modules already implement each piece.

\subsubsection{Deriving a hyperbolic system from the RTE: the P1
closure}
\label{sec:fuv-b-p1}

Take the first two angular moments of the radiative transfer equation
(the same starting point already used, correctly, in
Section~\ref{sec:fuv-gemini-review}'s critique of Gemini's derivation;
matching the notation of this codebase's own \code{rt/GEAR} module,
\texttt{theory/RadiativeTransfer/GEARRT/2-rt-equations.tex}):
\begin{align}
  \frac{\partial u}{\partial t} + \nabla\cdot\mathbf F
    &= -c\,\kappa_{\rm FUV}\rho\, u + S ,
  \label{eq:fuv-p1-zeroth} \\
  \frac{\partial \mathbf F}{\partial t} + c^2\,\nabla\cdot\mathbf P
    &= -c\,\kappa_{\rm FUV}\rho\, \mathbf F .
  \label{eq:fuv-p1-first}
\end{align}
$\mathbf P$ (the radiation pressure tensor) is unclosed -- it depends on
the second angular moment of the specific intensity, which these two
equations alone do not determine. The \textbf{Eddington closure},
\begin{equation}
  \mathbf P = \frac{u}{3}\,\mathbf I ,
  \label{eq:fuv-eddington-closure}
\end{equation}
is exact only when the local radiation field is close to isotropic
(frequent scattering/absorption randomising photon direction locally)
-- this closure, combined with Eqs.~\ref{eq:fuv-p1-zeroth}--\ref{eq:fuv-p1-first},
is the \textbf{P1 approximation}, a standard, well-studied closure of
the RTE (the ancestor of the M1 closure already implemented in this
fork's own \code{rt/GEAR} and \code{rt/SPHM1RT} modules,
Section~\ref{sec:fuv-b-precedent}), not something invented for this
logbook. Substituting Eq.~\ref{eq:fuv-eddington-closure} into
Eq.~\ref{eq:fuv-p1-first}:
\begin{equation}
  \frac{\partial \mathbf F}{\partial t} + \frac{c^2}{3}\,\nabla u
    = -\frac{\mathbf F}{\tau_{\rm phys}} , \qquad
  \tau_{\rm phys} = \frac{1}{c\,\kappa_{\rm FUV}\rho} .
  \label{eq:fuv-tau-phys}
\end{equation}
This is \textbf{exactly} the Cattaneo-type flux-relaxation equation
already integrated by \texttt{swiftsim\_0}'s
\code{hyperbolic/chemistry\_flux.h} (\code{dflux/dt + K/tau*grad Z =
-flux/tau}, with the exact analytic decay $\exp(-\Delta t/\tau)$
already used there for the source term -- the same
unconditional-stability property already exploited for the decay term
in Design~A, Section~\ref{sec:fuv-discretization}). Matching
Eq.~\ref{eq:fuv-tau-phys} to that form identifies
\begin{equation}
  D_{\rm phys} = c^2\,\tau_{\rm phys}/3 = \frac{c}{3\,\kappa_{\rm FUV}\rho}
   = \frac{c\,\lambda}{3} ,
  \qquad
  \tau_{\rm phys} = \frac{\lambda}{c} ,
  \label{eq:fuv-D-phys}
\end{equation}
the familiar radiative diffusion coefficient and photon mean-free-time
-- they fall out of the derivation here, rather than being asserted
independently as Gemini did. The hyperbolic signal speed of this exact
system is
\begin{equation}
  c_{\rm hyp,phys} = \sqrt{D_{\rm phys}/\tau_{\rm phys}} = \frac{c}{\sqrt3} ,
  \label{eq:fuv-chyp-phys}
\end{equation}
a \textbf{universal constant close to the speed of light}, independent
of opacity. \textbf{This is the catch}: taken at face value, the
rigorously RT-derived P1 system has exactly the same expensive
$\Delta t\sim h/c$ CFL that motivates reduced-speed-of-light
approximations in full M1 codes (\code{rt/GEAR}'s own
\code{reduced\_speed\_of\_light} parameter,
Section~\ref{sec:fuv-b-precedent}) -- P1 gives a real derivation, not
by itself a cheap one.

\subsubsection{Tying $\lambda$ to physics, $c_{\rm hyp}$ to
affordability: deriving $D$ and $\tau$}
\label{sec:fuv-b-lambda-chyp}

Approved 2026-07-24: rather than use the true, expensive
$(D_{\rm phys},\tau_{\rm phys},c_{\rm hyp,phys})$ of
Eqs.~\ref{eq:fuv-D-phys}--\ref{eq:fuv-chyp-phys}, decouple the two
physically-meaningful target quantities from the two free numerical
parameters. There are exactly two things that must be right, and two
free knobs ($D,\tau$) to set them with:
\begin{enumerate}
  \item the \textbf{screening length} $\lambda$
    (Eq.~\ref{eq:fuv-lambda-def}) must match real dust opacity -- this
    is the Yukawa/steady-state relation already derived in
    Section~\ref{sec:fuv-yukawa}, $\lambda=\sqrt{D\tau}$;
  \item the \textbf{propagation speed} $c_{\rm hyp}=\sqrt{D/\tau}$ is a
    cost/accuracy dial, chosen to be affordable rather than physically
    exact -- an honestly-named reduced-speed-of-light choice, the same
    kind of parameter \code{rt/GEAR} already exposes.
\end{enumerate}
Solving these two equations for the two unknowns,
\begin{equation}
  \tau = \frac{\lambda}{c_{\rm hyp}} , \qquad
  D = c_{\rm hyp}\,\lambda .
  \label{eq:fuv-lambda-chyp-solve}
\end{equation}
\textbf{$D$ is therefore not an independent free choice}: once
$\lambda$ is fixed by dust opacity (mandatory) and $c_{\rm hyp}$ is
picked for affordability (a genuine choice), $D$ and $\tau$ follow
algebraically. Concretely, choose
\begin{equation}
  c_{\rm hyp} = \kappa_{\rm margin}\,c_s ,
  \label{eq:fuv-chyp-choice}
\end{equation}
$c_s$ the local sound speed and $\kappa_{\rm margin}\sim5$--$10$ a
safety multiple (same margin-factor reasoning as the STARBENCH/Courant
discussion, Section~\ref{sec:fuv-b-cfl} below) -- large enough that the
field is quasi-instantaneous relative to hydro motions to a few
percent, without buying an unnecessarily tight (and unaffordable) CFL
by targeting anything close to $c_{\rm hyp,phys}=c/\sqrt3$.

\textbf{Why not use $D_{\rm phys}$/$\tau_{\rm phys}$ directly}: checked
against the ClumpSmith2021 regression-check snapshot (2026-07-23 run),
$\lambda=1/(\kappa_{\rm FUV}\rho)$ at solar metallicity comes out to
$\approx0.75\,$pc in the ambient gas ($n_H\approx130\,{\rm cm^{-3}}$,
$h\approx2.3\,$pc) and $\approx0.008\,$pc in the clump
($n_H\approx1.25\times10^4\,{\rm cm^{-3}}$, $h\approx0.5\,$pc) -- in
\emph{both} regions $\lambda<h$, meaning the "many scatterings per
resolution element" condition the diffusive regime
($\lambda_{\rm mfp}\ll h$) requires is not robustly satisfied even in
the diffuse gas at this resolution. The physical
$(D_{\rm phys},\tau_{\rm phys})$ pair is therefore not just
theoretically fragile in general (Section~\ref{sec:fuv-gemini-review})
but concretely not trustworthy in this project's own validated
configurations -- Eq.~\ref{eq:fuv-lambda-chyp-solve}'s decoupled choice
is the robust one, not a shortcut.

\subsubsection{Hyperbolic CFL}
\label{sec:fuv-b-cfl}

\begin{equation}
  \Delta t \le C_{\rm CFL}\,\frac{h}{c_{\rm hyp}} ,
  \label{eq:fuv-hyperbolic-cfl}
\end{equation}
\textbf{linear in $h$}, unlike Design~A's $h^2/D$ bound
(Section~\ref{sec:fuv-design-a-limit}). With $c_{\rm hyp}=\kappa_{\rm
margin}\,c_s$ (Eq.~\ref{eq:fuv-chyp-choice}) and the existing hydro
Courant condition $\Delta t_{\rm hydro}=C_{\rm CFL}\,h/c_s$, the ratio
\begin{equation}
  \frac{\Delta t_{\rm hyp}}{\Delta t_{\rm hydro}} = \frac{1}{\kappa_{\rm margin}}
  \label{eq:fuv-cfl-ratio}
\end{equation}
is independent of $h$ and $c_s$ entirely -- it only depends on the
chosen margin. With $\kappa_{\rm margin}\sim5$--$10$, this sits at or
below parity with the existing hydro step: essentially free, not the
prohibitive scaling Design~A's own $h^2$ bound would impose at the same
target speed (Section~\ref{sec:fuv-design-a-limit}).

\subsubsection{The concrete solve algorithm}
\label{sec:fuv-b-algorithm}

Putting Sections~\ref{sec:fuv-b-p1}--\ref{sec:fuv-b-cfl} together into
an actual per-timestep recipe, since "solve the hyperbolic system"
needs a concrete answer, not just the governing equations. State per
particle per band: $u_i$ (scalar) and $\mathbf F_i$ (3-vector) -- a
new field, not needed by Design~A.

\begin{enumerate}
  \item \textbf{Flux, per interacting pair.} Form a Rusanov flux
    (Eq.~\ref{eq:fuv-rusanov-flux}, Section~\ref{sec:fuv-b-rusanov})
    for the \emph{whole} state at once: the $u$-equation's flux is
    $\mathbf F\cdot\hat{\mathbf n}$ (already known -- no gradient needed
    to obtain it), the $\mathbf F$-equation's flux is $c^2\,\mathbf
    P\cdot\hat{\mathbf n}$ ($\mathbf P$ from the closure,
    Eq.~\ref{eq:fuv-eddington-closure} or~\ref{eq:fuv-m1-closure}), both
    using each particle's raw value at first order -- \textbf{no
    gradient is required for this step}; gradients
    (Section~\ref{sec:fuv-b-gradients}) are an optional upgrade to this
    step only, not a prerequisite for it.
  \item \textbf{Accumulate} pairwise, Newton's-third-law style, exactly
    like any symmetric SPH force loop.
  \item \textbf{Explicit hyperbolic sub-step}: update $(u,\mathbf F)$
    from the accumulated fluxes over $\Delta t$ -- the step that needs
    Eq.~\ref{eq:fuv-hyperbolic-cfl}.
  \item \textbf{Exact relaxation -- and this is where the screening
    term lands too.} Eq.~\ref{eq:fuv-p1-zeroth}'s own local sink,
    $-c\,\kappa_{\rm FUV}\rho\,u$ (the dust-screening/absorption term),
    is \emph{itself} a plain linear decay, $du/dt=-u/\tau_{\rm phys}$,
    with \textbf{exactly the same} $\tau_{\rm phys}=1/(c\,\kappa_{\rm
    FUV}\rho)$ (Eq.~\ref{eq:fuv-tau-phys}) already governing
    $\mathbf F$'s own Cattaneo relaxation -- not a coincidence, but a
    structural feature of the P1 system: both moments are being removed
    from the field by the \emph{same} physical absorption events, so
    they share the same rate. Both therefore admit \textbf{the same
    exact exponential integrator} already used for
    \texttt{swiftsim\_0}'s own hyperbolic module's flux relaxation
    (Section~\ref{sec:fuv-b-precedent}) and for Design~A's own decay
    term (Section~\ref{sec:fuv-discretization}):
    \begin{equation}
      u \to u\,e^{-\Delta t/\tau_{\rm phys}} , \qquad
      \mathbf F \to \mathbf F\,e^{-\Delta t/\tau_{\rm phys}} ,
      \label{eq:fuv-b-joint-relax}
    \end{equation}
    applied as one shared sub-step, unconditionally stable regardless
    of how $\Delta t$ compares to $\tau_{\rm phys}$.
  \item \textbf{Injection}: add the star's budget to $u$, as in
    Design~A.
\end{enumerate}

\textbf{This split is not optional, it is required by stiffness.}
Checked against the ClumpSmith2021 snapshot numbers already used in
Section~\ref{sec:fuv-b-lambda-chyp}: in the clump,
$\tau_{\rm phys}=\lambda/c_{\rm hyp}\approx0.008\,{\rm
pc}/(5\times0.33\,{\rm km\,s^{-1}})\approx5\times10^3\,$yr, shorter
than a typical hydro $\Delta t$ there ($\sim10^4$--$10^5\,$yr) --
$\Delta t\gg\tau_{\rm phys}$, a genuinely stiff relaxation. An explicit
(forward-Euler) treatment of Eq.~\ref{eq:fuv-b-joint-relax} would need
its own $\Delta t\ll\tau_{\rm phys}$ sub-cycling on top of the
hyperbolic CFL; the exact exponential integrator needs none, at zero
extra cost, for the identical reason Design~A's own decay term already
established (Section~\ref{sec:fuv-discretization}).

\textbf{Refinement available, not yet adopted}: the injection term
above is added \emph{after} the exact decay, as a separate additive
term -- accurate to $O(\Delta t)$ for the injection itself, exact only
for the pure decay. The fully exact solution of the inhomogeneous ODE
$du/dt=-u/\tau+S/m$ is $u(\Delta t)=u_0\,e^{-\Delta t/\tau}+(S\tau/m)
(1-e^{-\Delta t/\tau})$ -- folding injection into the same exact
integrator, rather than adding it afterwards, is a small, available
refinement to both designs, not yet applied to either.

\subsubsection{Gradients: an optional accuracy upgrade to the flux step}
\label{sec:fuv-b-gradients}

\textbf{Correction (2026-07-24) to an earlier draft of this section}:
the Brookshaw/Cleary-Monaghan pairwise operator (the tool actually used
by Design~A, Section~\ref{sec:fuv-discretization}) builds a
\emph{second} derivative ($\nabla^2u$) directly from paired \emph{value}
differences $(u_i-u_j)$ -- the right tool specifically when $\mathbf F$
has been eliminated and is not tracked as its own variable. Design~B
tracks $\mathbf F$ dynamically (Section~\ref{sec:fuv-b-algorithm}), so
it only ever needs \emph{first} derivatives -- $\nabla u$ (feeding
$\mathbf P$ via the closure) and $\nabla\cdot\mathbf F$ -- ordinary
SPH/kernel gradients, not the Brookshaw second-derivative machinery an
earlier pass of this section wrongly reused here.

\textbf{These gradients are optional, not required.} The baseline
algorithm (Section~\ref{sec:fuv-b-algorithm}, step~1) uses each
particle's raw $(u_i,\mathbf F_i)$ directly in the Rusanov flux -- a
first-order-in-space (piecewise-constant) scheme, complete on its own.
Gradients only enter if reconstructing $(u,\mathbf F)$ to the interface
for second-order spatial accuracy before forming the flux (a MUSCL-type
extrapolation) -- an accuracy upgrade on top of a working scheme, not a
prerequisite for it.

\textbf{Tier~1 vs Tier~2 for this optional step}: plain kernel gradient
$\nabla W_{ij}$ (Tier~1, matching this project's own
\code{GEAR\_DIFFUSION/chemistry\_iact.h}'s
\code{runner\_iact\_gradient\_diffusion}, no matrix inversion) versus
the matrix-corrected, linearly-consistent estimator using
$\mathbf B_i=\mathbf E_i^{-1}$ (Tier~2,
\texttt{swiftsim\_0}'s \texttt{GEAR\_FVPM\_Diffusion\_Theory.tex}'s own
gradient estimator).

\textbf{Correction (2026-07-24): $\mathbf B_i$ is not free in this
codebase's actual configuration.} Checked \code{config.h} for this
branch's build: \code{SPHENIX\_SPH} is defined, \emph{not}
\code{GIZMO\_MFM\_SPH}. \code{p->geometry.matrix\_E} is
GIZMO/MFM-specific infrastructure -- \texttt{swiftsim\_0}'s own
metal-diffusion module gets it "for free" only because that module
explicitly rides on a GIZMO/MFM hydro configuration (its own theory
doc states this outright). This project's production hydro is SPHENIX
(a density-based SPH scheme), which does not compute or maintain this
matrix at all. \textbf{Consequence}: Tier~2 here means standing up a
genuinely new per-particle least-squares moment-matrix computation
(its own kernel sum, its own $3\times3$ inversion) from scratch, not
reusing something the host hydro solver already provides -- a real,
previously uncounted engineering cost specific to this branch,
correcting the earlier "zero extra cost to access" claim. Tier~1
(plain kernel gradient, no matrix inversion) carries no such cost and
remains the practical default; Tier~2 should only be pursued if
Tier~1 (or the no-gradient baseline) fails a validation test.

\textbf{The symmetry subtlety, if Tier~2 is ever built}: $\mathbf B_i\ne
\mathbf B_j$ in general, so naively letting each particle use its own
$\mathbf B_i$ independently would break the antisymmetric pairwise
cancellation that keeps a symmetric exchange exactly conservative by
construction. A symmetrized average,
\begin{equation}
  \bar{\mathbf B}_{ij} = \tfrac12\left(\mathbf B_i + \mathbf B_j\right) ,
  \label{eq:fuv-tier15-bavg}
\end{equation}
used for both particles' shared link (the same harmonic-mean-family
discipline as Section~\ref{sec:fuv-harmonic-mean}, here an arithmetic
average of matrices rather than a harmonic mean of scalars, since it is
the shared gradient operator itself being built) fixes this, at the
same mild, not-yet-validated cost to the positivity/monotonicity
guarantee already flagged for Tier~2 gradients in general -- needing
the same empirical validation (Yukawa-profile and
intermediate-optical-depth tests, Section~\ref{sec:fuv-b-limits}) as
before this is trusted.

\subsubsection{Flux: Rusanov instead of a full Riemann solver}
\label{sec:fuv-b-rusanov}

A pure central-difference flux for a genuinely hyperbolic (wave)
equation is unconditionally unstable -- this is not specific to
radiative transfer, it is the standard reason upwind/Godunov methods
exist for any hyperbolic conservation law. Some dissipative treatment
is therefore unavoidable once Eq.~\ref{eq:fuv-tau-phys} is solved as a
genuinely time-dependent system rather than adiabatically eliminating
$\mathbf F$. \texttt{swiftsim\_0}'s own metal-diffusion module uses a
full Hopkins-2017-style HLL solver with wave-speed estimation and
minmod-blended artificial diffusivity (its own theory doc's flux
computation section) -- again, machinery justified there by
metal-mass conservation/positivity requirements this field does not
share.

The cheap alternative: a plain \textbf{Rusanov} (local Lax-Friedrichs)
flux \citep{Rusanov1961},
\begin{equation}
  F_{\rm face} = \tfrac12\left(F_L+F_R\right)
    - \tfrac12\,c_{\rm hyp}\left(U_R-U_L\right) ,
  \label{eq:fuv-rusanov-flux}
\end{equation}
one maximum-signal-speed, one dissipation term, no wave-speed
estimation subtlety, no minmod blending, no artificial-diffusivity
tuning -- technically the simplest member of the Riemann-solver family,
but a single line rather than the elaborate stack above. Valid because
$c_{\rm hyp}$ (Eq.~\ref{eq:fuv-chyp-choice}) is by construction the
system's own maximum characteristic speed. (\citet{Rusanov1961} is the
originating reference recalled from standard numerical-methods
literature -- author/year/general result are standard textbook
material, but the exact journal/volume/page details have not been
independently re-verified this pass, the same caveat already applied to
\citet{ClearyMonaghan1999} above.)

\subsubsection{Reusing the three hydro loops: task-graph cost}
\label{sec:fuv-b-tasks}

\textbf{Correction (2026-07-24) to an earlier draft's claim that this
design "needs its own new task subtype."} Checked directly against
this codebase's own task-graph machinery
(\code{src/runner\_doiact\_functions\_hydro.h} and this project's own
\code{src/chemistry/GEAR\_DIFFUSION/chemistry\_iact.h}): SWIFT's
existing hydro density/gradient/force pairwise loops
(\code{TASK\_LOOP\_DENSITY}, \code{TASK\_LOOP\_GRADIENT},
\code{TASK\_LOOP\_FORCE} branches of the same generated-function
template) are \emph{already} shared by multiple physics modules at
once -- e.g.\ \code{runner\_iact\_chemistry} is called inline inside
the density-loop branch, \code{runner\_iact\_gradient\_diffusion}
inside the gradient-loop branch, and \code{runner\_iact\_diffusion}
inside the force-loop branch, all within the same three generic
templates hydro itself uses. \textbf{No new task type or subtype is
needed}: adding a new physics channel to this codebase's existing
pattern means adding new \code{runner\_iact\_*} functions called from
these same three branches, not standing up new tasks in the scheduler
at all -- correcting the earlier, overly pessimistic claim.

\textbf{Concrete mapping for this design}:
\begin{itemize}
  \item \textbf{Density loop}: no new hook needed -- $\rho_i$ (feeding
    $\kappa_{\rm FUV}\rho$, Eq.~\ref{eq:fuv-lambda-def}) is already
    computed by hydro's own density loop regardless of this module;
  \item \textbf{Gradient loop}: a new \code{runner\_iact\_gradient\_isrf}
    (mirroring \code{runner\_iact\_gradient\_diffusion}) -- \emph{only}
    needed if the optional gradient upgrade of
    Section~\ref{sec:fuv-b-gradients} is adopted; the baseline
    algorithm (Section~\ref{sec:fuv-b-algorithm}) does not need this
    loop hook at all;
  \item \textbf{Force loop}: a new \code{runner\_iact\_isrf}/
    \code{runner\_iact\_nonsym\_isrf} (mirroring
    \code{runner\_iact\_diffusion}/\code{runner\_iact\_nonsym\_diffusion})
    -- the one hook the baseline algorithm actually requires, computing
    the Rusanov flux (Eq.~\ref{eq:fuv-rusanov-flux}) and accumulating
    the explicit hyperbolic sub-step
    (Section~\ref{sec:fuv-b-algorithm}, steps~1--3).
\end{itemize}
\textbf{Minimal buildable version}: one new force-loop hook only, no
gradient loop, no new task type -- cheaper, in task-graph terms, than
the earlier draft credited it for. This does not exempt the new field
($\mathbf F_i$, Section~\ref{sec:fuv-b-algorithm}) from the usual
cell-link/\code{cell\_clean\_links} bookkeeping this project's own
task-registration checklist requires for any new per-particle state,
nor from choosing which existing timestep/rebuild cadence machinery
(Section~\ref{sec:fuv-discretization}'s adaptive-rebuild precedent)
governs how often it updates.

\textbf{Design~A, for comparison}: its single kernel-average update
(Eq.~\ref{eq:fuv-discrete-update}) needs only the density-loop-style
hook (matching plain \code{runner\_iact\_chemistry}'s own one-loop
pattern) -- no gradient loop, no force loop, since there is no
separate gradient reconstruction or flux exchange step in Design~A at
all. Both designs turn out to need no new task types; they differ in
\emph{how many} of the three existing loop hooks they use (one for
Design~A; one to two for Design~B, depending on whether the optional
gradient upgrade is adopted), not in whether new scheduler machinery is
needed at all.

\subsubsection{From P1 to a cheap M1: a local closure upgrade}
\label{sec:fuv-b-m1}

\textbf{The difference between P1 and M1.} P1's closure
(Eq.~\ref{eq:fuv-eddington-closure}) is \emph{always} $\mathbf P=u/3\,
\mathbf I$, regardless of how directional the true local field actually
is -- exact only near isotropy. The \textbf{M1 closure}
\citep{Levermore1984} -- already implemented in this exact SWIFT fork,
\code{rt/GEAR} and \code{rt/SPHM1RT}, and transcribed verbatim from
\texttt{theory/RadiativeTransfer/GEARRT/2-rt-equations.tex} -- instead
makes the Eddington tensor a function of the local \textbf{reduced
flux} $f=|\mathbf F|/(cu)$:
\begin{align}
  \mathbf P &= \mathbf D\,u , \qquad
  \mathbf D = \frac{1-\chi}{2}\,\mathbf I + \frac{3\chi-1}{2}\,
    \hat{\mathbf n}\otimes\hat{\mathbf n} , \qquad
  \hat{\mathbf n}=\frac{\mathbf F}{|\mathbf F|} ,
  \label{eq:fuv-m1-closure} \\
  \chi(f) &= \frac{3+4f^2}{5+2\sqrt{4-3f^2}} .
  \label{eq:fuv-m1-chi}
\end{align}
Check the limits: $\chi(0)=1/3$ (recovers P1 exactly -- correct
isotropic/diffusive behaviour) and $\chi(1)=1$ (correct free-streaming
behaviour -- pure advection along $\hat{\mathbf n}$, no unphysical
sideways spreading). \textbf{This is why M1 is better}: it is not
locked into the isotropic assumption everywhere, so near a strong,
little-scattered point source ($f\to1$) it lets the field stream
directionally instead of leaking sideways the way P1 unconditionally
forces it to -- the mechanism that gets shadows mostly right where P1
cannot.

\textbf{Could we upgrade to a cheap M1? Yes.} The physics upgrade and
the numerics can be kept separate, exactly as in the gradient
discussion above:
\begin{itemize}
  \item \textbf{Cheap part}: the closure swap itself needs no new
    state -- $(u,\mathbf F)$ is already tracked by the hyperbolic system
    of Section~\ref{sec:fuv-b-p1}, and $f$, $\chi(f)$, $\mathbf D$ are
    local, per-particle/per-interface formula evaluations, essentially
    free next to the cost of a neighbour loop.
  \item \textbf{Expensive part, deliberately not adopted}: the
    \emph{proper} M1 solver already in \code{rt/GEAR}
    (\code{rt\_riemann\_HLL.h}) needs a precomputed
    $100\times100$ lookup table of eigenvalues
    \citep{Gonzalez2007}, because M1's true characteristic speeds are a
    nontrivial function of both $f$ and the angle between the flux and
    the interface normal -- genuinely more complex than P1's simple
    $\pm c_{\rm hyp}$. Reusing the cheap Rusanov flux of
    Eq.~\ref{eq:fuv-rusanov-flux} instead remains valid, because M1's
    true characteristic speeds are bounded within $[-c_{\rm hyp},
    c_{\rm hyp}]$ by construction (M1's defining design property is that
    the reduced flux never exceeds 1) -- a Rusanov flux using
    $c_{\rm hyp}$ as the dissipation speed is a correct, stable bound
    for the upgraded system too, just not as sharp (less numerically
    diffusive) as the proper eigenvalue-based solver would be.
\end{itemize}
So: keep M1's closure (Eqs.~\ref{eq:fuv-m1-closure}--\ref{eq:fuv-m1-chi}),
keep the cheap Rusanov flux (Eq.~\ref{eq:fuv-rusanov-flux}) -- most of
M1's physical upgrade over P1, essentially none of its own numerical
machinery cost.

\textbf{Implementation subtlety to get right}: under a
reduced-speed-of-light scheme, $f$ must be built with the \emph{same}
reduced $c_{\rm hyp}$, not the true $c$. If $f=|\mathbf F|/(c\,u)$ used
true $c$ while $\mathbf F$ is dynamically bounded by $c_{\rm hyp}\,u$
(with $c_{\rm hyp}\ll c$), $f$ would always come out tiny regardless of
how beamed the field really is, silently collapsing the upgraded
closure back to P1-like behaviour and defeating the point of upgrading.
(General RT-methods knowledge, not independently verified against
\code{rt/GEAR}'s own \code{reduced\_speed\_of\_light} handling this
pass -- worth checking there too, not only in this design, before
either is trusted.)

\subsubsection{Known, accepted limitations: shadows and crossing beams}
\label{sec:fuv-b-limits}

\textbf{P1's shadow-smearing} (already accepted as a trade-off,
Section~\ref{sec:fuv-b-p1}'s validity caveat): a real photon field shows
almost zero flux directly behind a dense clump and full flux to the
side; a diffusion-type (P1) field leaks some flux around the clump via
neighbouring particles, smearing that shadow into something softer than
reality. The cheap-M1 upgrade above improves this (Eq.~\ref{eq:fuv-m1-chi}
lets $f\to1$ near a strong, little-scattered source) but does not fully
fix it.

\textbf{M1's own, distinct limitation: crossing beams.} M1 is a
single-direction closure -- at any point it assumes all local flux
points one way, $\hat{\mathbf n}=\mathbf F/|\mathbf F|$. When two
genuinely separate beams from two different sources cross at a gas
particle, M1 cannot represent "two beams passing through each other
unaffected" -- it merges them into one averaged direction, a
spurious artifact at the crossing point, well documented in the RT
literature and not fixed by a better (more expensive) M1 solver; it is
intrinsic to the closure. Directly relevant here since this project
already cares about multi-source/angular behaviour (the HEALPix HII
splitting, Section~\ref{sec:algorithm}).

\textbf{Why, physically: moment-closure RT is to full kinetic RT what
fluid dynamics is to kinetic gas dynamics} (Darwin's framing,
2026-07-24, confirmed correct). Both P1/M1 and ordinary Euler
hydrodynamics are moment truncations of an underlying kinetic
(Boltzmann-type) transport equation -- radiative transfer's RTE for
photons, the Boltzmann equation for a gas of particles. Both inherit
the identical failure mode: a moment/fluid description collapses any
locally multi-modal distribution (two distinct particle streams that
have not thermalised together; two distinct photon beams that have not
scattered together) into one local bulk description (one flow velocity;
one flux direction), because a finite set of angular moments simply
cannot distinguish "two populations moving differently" from "one
population with the combined average motion." Representing genuinely
crossing, non-interacting populations correctly requires a kinetic (or
here, ray-based/multi-group) description, not a better moment closure --
this is a structural property of the moment-truncation approach, not a
correctable numerical shortcoming.

\subsubsection{Relation to existing SWIFT modules}
\label{sec:fuv-b-precedent}

This design borrows validated formulas and structure from three places
already in this SWIFT fork rather than inventing them:
\begin{itemize}
  \item \texttt{swiftsim\_0}'s \code{chemistry/GEAR\_FVPM\_DIFFUSION/hyperbolic/}
    -- the Cattaneo relaxation-equation structure
    ($\partial\mathbf F/\partial t + (K/\tau)\nabla u=-\mathbf F/\tau$,
    $c_{\rm hyp}=\sqrt{\|K\|/\tau}$), its exact analytic
    ($\exp(-\Delta t/\tau)$) integration of the relaxation source term,
    and its hyperbolic CFL ($\Delta t=C_{\rm CFL}\,\Delta x/v_{\max}$) --
    all mirrored directly in Sections~\ref{sec:fuv-b-p1}--\ref{sec:fuv-b-cfl}
    above, for a different transported quantity;
  \item \texttt{swiftsim\_0}'s \code{chemistry/GEAR\_FVPM\_DIFFUSION/}
    (parabolic) gradient estimator, \code{p->geometry.matrix\_E}/
    $\mathbf B_i=\mathbf E_i^{-1}$ -- the source of the optional Tier~2
    correction matrix, Section~\ref{sec:fuv-b-gradients} (not free in
    this project's own SPHENIX-based configuration, unlike there);
  \item this project's own \code{chemistry/GEAR\_DIFFUSION/chemistry\_iact.h}
    and \code{src/runner\_doiact\_functions\_hydro.h} -- confirms new
    physics can hook into the existing density/gradient/force loops via
    new \code{runner\_iact\_*} functions, with no new task type,
    Section~\ref{sec:fuv-b-tasks};
  \item \code{rt/GEAR} and \code{rt/SPHM1RT} -- full, working M1-closure
    RT solvers already in this codebase, confirming
    Eqs.~\ref{eq:fuv-m1-closure}--\ref{eq:fuv-m1-chi} exactly (verbatim
    match against \texttt{theory/RadiativeTransfer/GEARRT/2-rt-equations.tex}),
    and the source of the \code{reduced\_speed\_of\_light} precedent
    (\code{rt/GEAR/rt\_properties.h}) for Eq.~\ref{eq:fuv-chyp-choice},
    plus the Gonzalez-2007 eigenvalue-table solver
    (\code{rt/GEAR/rt\_riemann\_HLL.h}) this design deliberately avoids
    reusing directly (Section~\ref{sec:fuv-b-m1}) in favour of the
    cheaper Rusanov flux.
\end{itemize}
None of these are proposed for direct reuse as full modules -- they are
full, more expensive solvers built for different purposes (exact metal
mass conservation; full radiation hydrodynamics) -- but each already
validates a specific piece of this design's own physics or numerics,
which is worth more than an independent derivation would be.

\subsubsection{Open items, Design B}
\label{sec:fuv-design-b-open}

\begin{itemize}
  \item \textbf{Resolved (2026-07-24)}: task-graph placement is no
    longer an open question -- Section~\ref{sec:fuv-b-tasks} confirms
    new \code{runner\_iact\_*} hooks into the existing density/
    gradient/force loops, no new task type, mirroring
    \code{GEAR\_DIFFUSION}'s own precedent directly. What remains open:
    building the optional Tier~2 gradient upgrade
    (Section~\ref{sec:fuv-b-gradients}) needs a genuinely new
    per-particle moment-matrix computation in this SPHENIX-based
    codebase (not free, unlike \texttt{swiftsim\_0}'s MFM-based
    checkout), and the mild positivity/monotonicity loosening it
    introduces is not yet empirically validated;
  \item \citet{Rusanov1961}'s exact bibliographic details are not
    independently re-verified;
  \item the cheap-M1 upgrade's requirement that $f$ use the reduced
    $c_{\rm hyp}$, not true $c$ (Section~\ref{sec:fuv-b-m1}), is not yet
    implemented or checked against \code{rt/GEAR}'s own handling of the
    same issue;
  \item the margin factor $\kappa_{\rm margin}$
    (Eq.~\ref{eq:fuv-chyp-choice}) is proposed at $\sim5$--$10$ by
    analogy with a general "quasi-instantaneous relative to hydro"
    argument, not yet empirically tuned or validated against an actual
    run;
  \item folding the injection term into the same exact exponential
    integrator as the joint $u$/$\mathbf F$ relaxation
    (Section~\ref{sec:fuv-b-algorithm}'s closing remark) is an available
    refinement, not yet applied;
  \item a validation plan specific to this design is not yet built:
    the Yukawa-profile and telegrapher-wavefront tests already proposed
    for Design~A apply equally here, plus a new
    intermediate-optical-depth/off-axis-clump ("partial shadow") test
    to quantify the known P1/M1 shadowing limitation
    (Section~\ref{sec:fuv-b-limits}) for this project's own realistic
    clump opacities, rather than asserting it away;
  \item whether Design~A or Design~B is actually implemented first --
    or whether Design~A is kept as a cheap fallback/cross-check target
    for Design~B's own quasi-steady ($\tau\to0$-like) limit -- is not
    yet decided;
  \item item~T6 as a whole remains unscoped as an implementation, for
    either design.
\end{itemize}

\subsection{Open items common to both designs}
\label{sec:fuv-open-common}

\begin{itemize}
  \item confirming whether \citet{DraineBertoldi1996}'s $k_{\rm
    LW}(\chi)$ normalization, or a direct $\sigma_{\rm LW}$-based
    conversion (Section~\ref{sec:fuv-two-bands}), is the better fit
    once $L_{\rm LW}$ is tracked explicitly;
  \item checking whether \code{use\_isrf\_field} and
    \code{use\_radiative\_transfer} can be enabled simultaneously
    without conflict (not checked -- they gate physically distinct
    code paths in \code{cool1d\_multi\_g.F} above, but this has not
    been tested end to end).
\end{itemize}
